
%%
%% Test: matinbook.cls — Complete Class Bootstrap
%% Stage: v1.1 (stages 1-4)
%% 
%% Purpose: Verify all fixes applied to `matinbook.cls` in v1.1:
%%   Stage 1 (LaTeX 2023 standards):
%%     - No relative paths in \RequirePackage (uses \input instead)
%%     - \NeedsTeXFormat has a release date
%%     - \ClassInfo used instead of \typeout
%%     - No obsolete "minimal theme" reference
%%     - Copyright header updated (2025-2026) + repository URL
%%   Stage 2 (Official documentation):
%%     - \ProvidesClass reports v1.1
%%   Stage 3 (Architecture patterns):
%%     - Bootstrap sequence documented in header
%%     - Structured section comments
%%     - Welcome message with accurate module counts
%%   Stage 4 (Market standards):
%%     - Paper size documented (A4 now, vaziri deferred to mb-layout)
%% 

\documentclass{matinbook}

%----------------------------------------
% Test Metadata
%----------------------------------------
\title{تست کلاس MatinBook — نسخه ۱.۱}
\author{تیم MatinBook}
\date{۱۴۰۵}

\begin{document}

\maketitle

\chapter{بررسی اصلاحات نسخه ۱.۱}
\label{chap:test-cls}

این سند یک تست جامع برای بررسی اصلاحات اعمال‌شده در فایل
\verb|matinbook.cls| نسخه ۱.۱ است. هر بخش، یک یا چند مورد از
مراحل چهارگانه‌ی بازبینی را پوشش می‌دهد.

%----------------------------------------
\section{مرحله ۱ — استانداردهای LaTeX 2023}
\label{sec:stage1}

\subsection{رفع مسیرهای نسبی}

در نسخه‌ی قبل، ماژول‌ها با \verb|\RequirePackage{tex/...}| بارگذاری
می‌شدند که از نظر استاندارد LaTeX نادرست است. در نسخه‌ی ۱.۱، همه‌ی
ماژول‌ها با \verb|\input{tex/...sty}| بارگذاری می‌شوند.

\textbf{نتیجه:} اگر این سند بدون خطا کامپایل شود، یعنی مسیرها
به‌درستی resolve شده‌اند.

\subsection{تاریخ در \texttt{\textbackslash NeedsTeXFormat}}

نسخه‌ی جدید از \verb|\NeedsTeXFormat{LaTeX2e}[2023/06/01]| استفاده
می‌کند که استاندارد LaTeX 2023 است.

\subsection{استفاده از \texttt{\textbackslash ClassInfo}}

پیام خوش‌آمد حالا با \verb|\ClassInfo{matinbook}{...}| چاپ می‌شود،
نه \verb|\typeout|. این یعنی در حالت \verb|\batchmode| ترمینال
شلوغ نمی‌شود.

\textbf{بررسی:} به فایل \verb|stage-cls-01.log| مراجعه کنید و
خطوط شروع‌شده با \verb|Class matinbook Info:| را ببینید.

%----------------------------------------
\section{مرحله ۲ — مستندات رسمی}
\label{sec:stage2}

\subsection{نسخه‌ی کلاس}

\verb|\ProvidesClass| حالا نسخه‌ی \verb|v1.1| را گزارش می‌دهد
(قبلاً \verb|v1.0| بود).

\textbf{بررسی:} به لاگ کامپایل نگاه کنید و بدنبال
\verb|matinbook.cls ... v1.1| بگردید.

%----------------------------------------
\section{مرحله ۳ — الگوهای معماری}
\label{sec:stage3}

\subsection{ساختار مستندسازی}

هدر فایل حالا شامل جدول \textbf{BOOTSTRAP SEQUENCE} است که ترتیب
بارگذاری ماژول‌ها را توضیح می‌دهد. این الگو از \verb|classicthesis|
گرفته شده است.

\subsection{پیام Welcome با شمارش دقیق}

شمارش ماژول‌ها در پیام Welcome اصلاح شده:

\begin{itemize}
    \item Core: ۴ ماژول
    \item Typography: ۳ ماژول
    \item Locale: ۱ ماژول
    \item Features: ۸ ماژول
    \item Systems: ۳ ماژول
    \item Theme: ۱ سیستم
    \item \textbf{مجموع:} ۱۹ ماژول + سیستم تم
\end{itemize}

%----------------------------------------
\section{مرحله ۴ — استانداردهای بازار}
\label{sec:stage4}

\subsection{قطع کتاب}

قطع فعلی \verb|a4paper| است (برای سازگاری با نسخه‌ی قبل).
قطع وزیری (۱۶.۵×۲۳.۵ سانتی‌متر) در \verb|mb-layout.sty| پیاده‌سازی
خواهد شد.

\textbf{بررسی ابعاد صفحه:}

\begin{quote}
\verb|\the\paperwidth| = \the\paperwidth \\
\verb|\the\paperheight| = \the\paperheight
\end{quote}

اگر مقدار \verb|\the\paperwidth| برابر \verb|597.50787pt| باشد
(یعنی ۲۱ سانتی‌متر)، قطع A4 تأیید می‌شود.

%----------------------------------------
\section{محیط‌های علمی}
\label{sec:environments}

برای اطمینان از اینکه همه‌ی ماژول‌ها به‌درستی بارگذاری شده‌اند،
چند محیط نمونه را آزمایش می‌کنیم:

\subsection{قضیه}

\begin{theorem}[قضیه‌ی نمونه]
    برای هر عدد حقیقی $x$، داریم $x^2 \geq 0$.
\end{theorem}

\begin{proof}
    چون مربع هر عدد حقیقی نامنفی است، نتیجه حاصل می‌شود.
\end{proof}

\subsection{تعریف}

\begin{definition}[عدد اول]
    عدد طبیعی بزرگ‌تر از ۱ که تنها دو مقسوم‌علیه مثبت دارد،
    \emph{عدد اول} نامیده می‌شود.
\end{definition}

\subsection{مثال}

\begin{example}[اعداد اول کوچک]
    اعداد $2, 3, 5, 7, 11, 13$ نمونه‌هایی از اعداد اول هستند.
\end{example}

%----------------------------------------
\section{ریاضیات}
\label{sec:math}

بررسی نمادهای ریاضی سفارشی:

\begin{equation}
    \label{eq:test}
    f(x) = \sum_{i=1}^{n} x_i^2 + \int_{0}^{\infty} e^{-t^2} \, dt
\end{equation}

مجموعه‌های اعداد: $\N \subset \Z \subset \Q \subset \R \subset \C$.

قدر مطلق: $\abs{-5} = 5$، نُرم: $\norm{\vec{v}}$،
کف و سقف: $\floor{3.7} = 3$، $\ceil{3.2} = 4$.

%----------------------------------------
\section{کد}
\label{sec:code}

\begin{latin}
    \begin{minted}{python}
    def factorial(n):
        """Compute n! recursively."""
        if n <= 1:
            return 1
        return n * factorial(n - 1)
    \end{minted}
\end{latin}

%----------------------------------------
\section{نتیجه‌گیری}
\label{sec:conclusion}

اگر این سند بدون خطا کامپایل شده باشد، یعنی:

\begin{itemize}
    \item ✅ مسیرهای \verb|\input| به‌درستی resolve شده‌اند
    \item ✅ همه‌ی ۱۹ ماژول بارگذاری شده‌اند
    \item ✅ نسخه‌ی \verb|v1.1| فعال است
    \item ✅ پیام Welcome در لاگ ثبت شده است
    \item ✅ محیط‌های علمی و ریاضی کار می‌کنند
    \item ✅ سیستم کد (minted) فعال است
\end{itemize}

\textbf{وضعیت تست:} ✅ قبول

\end{document}
